\chapter{Reproducibility}
\label{app:reproducibility}

\section{Software and random seeds}

The primary experiment was run with Python 3.13.5, NumPy 2.4.4, SciPy
1.17.1, and pandas 3.0.3. Its population-generation seed was 2,026,080,501
and its simulation seed was 2,026,080,502. The scaled chi-squared audit used
seeds 2,026,080,701 and 2,026,080,702. The exact settings and software
versions are stored beside each output in \texttt{run\_metadata.json}.

\section{Primary experiment}

From the repository root, the primary experiment is reproduced with
\begin{verbatim}
MPLBACKEND=Agg MPLCONFIGDIR=$PWD/.mplcache \
  .venv/bin/python \
  WelchSatterthwaiteMI/experiments/run_supervisor_experiment.py \
  --profile full \
  --output-dir WelchSatterthwaiteMI/results/supervisor_practical
\end{verbatim}
The full profile uses 10,000 null replicates per scenario, 10,000 replicates
per power setting, batches of 1,000, and 200 timing repetitions. The recorded
SHA-256 hash of the experiment script was
\begin{verbatim}
9e654a849ff100eedd57c2703ee5617a32d7c35b3400ef37be718adca35f9123
\end{verbatim}

The output directory contains the following principal artefacts:
\begin{itemize}[leftmargin=*]
  \item \texttt{population\_scenarios.csv}:\\
  Fixed population probabilities, sample sizes, generating parameters, and
  expected-count diagnostics for all 60 null scenarios.

  \item \texttt{scenario\_results.csv}:\\
  Scenario-level FPRs, calibration errors, Wilson intervals, coverage, and
  validity diagnostics.

  \item \texttt{regime\_summary.csv}:\\
  Equal-weight summaries over the 12 population pairs in each regime.

  \item \texttt{null\_pvalues.npz}:\\
  Replicate-level p-values used to construct the rejection-calibration curves.

  \item Calibration-grid files:\\
  \texttt{rejection\_calibration\_scenarios.csv} and\\
  \texttt{rejection\_calibration\_regimes.csv}.\\
  Actual rejection rates over the complete nominal-alpha grid.

  \item \texttt{power\_summary.csv} and \texttt{runtime\_summary.csv}:\\
  Alternative-hypothesis and computational results.
\end{itemize}

\section{Scaled chi-squared audit}

The component audit is reproduced with
\begin{verbatim}
MPLBACKEND=Agg MPLCONFIGDIR=$PWD/.mplcache \
  .venv/bin/python \
  WelchSatterthwaiteMI/experiments/validate_scaled_chi_square.py \
  --profile focused \
  --replicates 10000 \
  --output-dir \
  WelchSatterthwaiteMI/results/scaled_chi_square_validation
\end{verbatim}
The corresponding output directory contains component-, regime-, shape-, and
model-level summaries together with the fixed populations and metadata.

\section{Automated checks}

The implementation tests are run with
\begin{verbatim}
.venv/bin/python -m unittest discover \
  -s WelchSatterthwaiteMI/tests -v
\end{verbatim}
The tests cover MI calculation, bias correction, normal and Welch references,
the variance influence calculation, table validation, reproducibility of the
experiment profiles, and scaled chi-squared audit utilities.

\section{Building this document}

The scenario-calibration, effective-degrees-of-freedom, and power figures are
generated from the accepted CSV outputs. From
\texttt{Thesis Writeup/Welch MI}, run
\begin{verbatim}
MPLBACKEND=Agg MPLCONFIGDIR=/tmp/welch-mi-mpl \
  ../../.venv/bin/python figures/make_thesis_figures.py
\end{verbatim}
The thesis is then built with
\begin{verbatim}
/Library/TeX/texbin/latexmk -pdf \
  -interaction=nonstopmode -halt-on-error main.tex
\end{verbatim}
The bibliography is read from
\texttt{WelchSatterthwaiteMI/article/references.bib}. The calibration figure
and all thesis-specific summary figures are generated from the accepted
experiment outputs rather than from separately edited values.
